\section{Weather Bottlenecks: A Distinct Type}

ATRI's congestion rankings capture recurring traffic delay. They do not capture weather-driven closures, which impose a fundamentally different economic cost pattern: low frequency, long duration, full closure (capacity = 0), concentrated on specific mountain pass segments.

\subsection{Donner Pass (I-80)}

Interstate 80 crosses the Sierra Nevada at Donner Pass (elevation 7,227 feet) approximately 90 miles east of Sacramento. The pass closes approximately 50 times per year for weather, with a mean closure duration of 18 hours \citep{Caltrans2023}. The seasonal pattern is concentrated in November-April; the worst months (December-February) can see 10+ closures.

When Donner closes, eastbound I-80 freight has three options:
\begin{enumerate}
\item Wait for reopening (mean 18h, maximum 72h)
\item Reroute via I-40 through southern California (+310 miles, +5.5 hours driving)
\item Reroute via US-50 through Lake Tahoe (not maintained for winter truck traffic; not viable for most commercial vehicles)
\end{enumerate}

Annual cost estimate: 50 closures $\times$ 18h mean duration $\times$ 8,000 trucks/day at Donner $\times$ \$225/hr rerouting premium = \textbf{\$1.6 billion/year}.

This cost is approximately equal to ATRI bottleneck locations 8--10 combined (\$1.98B for I-75 Atlanta, I-95 Baltimore, I-80 Chicago). Donner Pass is the 11th most expensive freight bottleneck in the US if weather closures are included in the ranking --- and it does not appear in ATRI's top 100 because ATRI measures recurring congestion, not episodic closures.

\subsection{Snoqualmie Pass (I-90)}

Interstate 90 crosses the Cascade Mountains at Snoqualmie Pass (elevation 3,022 feet), approximately 60 miles east of Seattle. The pass closes approximately 40 times per year, with a mean duration of 12 hours \citep{WSDOT2023}. Unlike Donner, Snoqualmie has a viable alternate: US-2 over Stevens Pass (lower volume, adds 60 miles). But US-2 is not maintained for commercial truck traffic during winter closures.

Annual cost estimate: 40 closures $\times$ 12h $\times$ 6,000 trucks/day $\times$ \$225/hr = \textbf{\$0.65 billion/year}.

\subsection{Implications for I2.0}

Weather bottlenecks require a fundamentally different intervention than congestion bottlenecks:
\begin{itemize}
\item Congestion bottlenecks: managed freight lanes to add capacity during peak periods
\item Weather bottlenecks: new infrastructure to bypass or protect the exposed segment (tunnel, alternate alignment, snowshed, or accelerated clearance technology)
\end{itemize}

The I2.0 investment case for Donner Pass is a \$4B freight tunnel that eliminates weather closure risk entirely --- a 2.5-year payback period at \$1.6B/year in avoided costs. No other single infrastructure investment in the I2.0 program has this NPV per dollar of construction cost.
